\documentclass[11pt,twocolumn]{article}

\usepackage[margin=0.6in]{geometry}
\usepackage[utf8]{inputenc}
\usepackage[T1]{fontenc}
\usepackage{ebgaramond}

\usepackage{bm}
\newcommand{\unit}[1]{\hat{\bm {#1}}}
\newcommand{\parens}[1]{\left(#1\right)}

\title{Astranesse}
\author{Jack T. Dinsmore}

\renewcommand{\baselinestretch}{1.15}

\begin{document}
  \maketitle

  \section{Planet generation}
  \subsection{Atmosphere}
    Given the basic properties of a planet and a rough altitude map the task is to compute its atmospheric flows. These will then be used to determine the weathering and humidity maps of the planet, which modify the altitude and set the biome. The atmospheric simulation should include Hadley cells and respond to large scale features such as coastlines and mountains. The humidity can be done post hoc, and should take adiabatic cooling into account.

    My basic picture is that Hadley cells exist because the higher temperature at the equator sets up a pressure gradient towards the poles. But the upper atmosphere has far less pressure gradient because it's less thermally connected to the ground. This sets up curl in the pressure field.

  \subsection{Fog}
    Fog consists of the combined effects of scattering out of the line of sight and into the line of sight. Starting with the former (extinction),

    \subsubsection{Extinction}
      I'll use the optical depth law
      \begin{equation}
        \tau = \int_0^\infty d\ell\, E(\ell) = \int_0^\infty d\ell\, E_0 e^{-h(\ell) / H}
      \end{equation}
      where $E(\ell)$ is the differential extinction and the fog color is blended into the background color by the amount $1 - e^-\tau$. The second equation comes from assuming an exponentially decaying atmosphere with surface extinction $E_0$ (inverse meters), scale height $H$, and height $h(\ell)$ above the surface.

      For a cubic planet, $h(\ell)$ depends on what face the point of integration is hovering over. So I'll split the trajectory into chunks, where each chunk is over a unique edge, face, or vertex. One way of formalizing this is to define the vector $\bm \ell$ pointing from the camera to the end of the sight ray, and dividing the vector into chunks ranging from $\alpha_i \bm \ell$ to $\alpha_{i+1} \bm \ell$, so that
      \begin{equation}
        \tau = |\ell| E_0 \sum_i \int_{\alpha_i}^{\alpha_{i+1}}d\alpha\, e^{-h_i(\alpha) / H}.
      \end{equation}

      If the chunk is over a face with normal vector $\unit n$, then the height is $(\bm a + \alpha \bm \ell) \cdot \unit n - w$ where $w$ is the cube half-width and $\bm a$ is the ray's starting position. The integral of the summand is then
      \begin{equation}
        \tau_i = E_0H\frac{\ell}{\bm \ell \cdot \unit n_i} e^{-(\bm a \cdot \unit n_i - w_i) / H} \parens{e^{-\bm\ell \cdot \unit n_i \alpha_{i} / H} - e^{-\unit\ell \cdot \unit n_i \alpha_{i+1} / H}}.
      \end{equation}
      It seems that one easy way to continue this formula for edges and corners is just to take the normal vector as the average of the abutting faces' normal vectors.

      Note that $E_0H$ is dimensionless. I call that the extinction constant. Then the coordinate system can be remapped in whatever way is desired, e.g.~to a frame in which the planet is the unit cube.

\end{document}